\documentclass[10pt,a4paper]{article}
\usepackage{amsmath,amssymb,amsthm,graphicx}
\usepackage{hyperref}
\usepackage{mathtools}
\usepackage{nicefrac}
\usepackage{xspace}
\usepackage{xcolor}
\usepackage{url}
\usepackage{bm}
\usepackage{enumitem}
\newcommand{\mathbbm}[1]{\mathbf{#1}}
\newtheorem{theorem}{Theorem}
\newtheorem{thm}{Theorem}
\newtheorem{conj}[thm]{Hypothesis}
\newtheorem{rconj}[thm]{Conjecture}
\newtheorem{lemma}[thm]{Lemma}
\newtheorem{corollary}[thm]{Corollary}
\newtheorem{claim}[thm]{Claim}
\newtheorem{proposition}[thm]{Proposition}
\newtheorem{definition}[thm]{Definition}
\newtheorem{remark}[thm]{Remark}
\newtheorem{spec}[thm]{Speculation}
\newtheorem{fact}[thm]{Fact}
\newtheorem{problem}{Problem}
\newcommand\comment[1]{\fbox{\tt #1}}
\newcommand\E{\mathop{\mathbb{E}}}
\newcommand\Var{\mathop{Var}}

\newcommand\card[1]{\left| {#1} \right|}
\newcommand\sett[2]{\left\{ \left. #1 \;\right\vert #2 \right\}}
\newcommand\settright[2]{\left\{ #1 \;\left\vert\; #2 \right.\right\}}
\newcommand\set[1]{{\left\{ #1 \right\}}}
\newcommand\Prob[2]{{\Pr_{#1}\left[ {#2} \right]}}
\newcommand\BProb[2]{{\Pr_{#1}\Big[ {#2} \Big]}}
\newcommand\bProb[2]{{\Pr_{#1}\big[ {#2} \big]}}
\newcommand\cProb[3]{{\Pr_{#1}\left[ \left. #3 \;\right\vert #2 \right]}}
\newcommand\cbProb[3]{{\Pr_{#1}\big[ \big. #3 \;\big\vert #2 \big]}}
\newcommand\cBProb[3]{{\Pr_{#1}\Big[ \Big. #3 \;\Big\vert #2 \Big]}}
\newcommand\Expect[2]{{\mathop{\mathbb{E}}_{#1}\left[ {#2} \right]}}
\newcommand\bExpect[2]{{\mathop{\mathbb{E}}_{#1}\big[ {#2} \big]}}
\newcommand\BExpect[2]{{\mathop{\mathbb{E}}_{#1}\Big[ {#2} \Big]}}
\newcommand\cProbright[3]{{\Pr_{#1}\left[ #3\;\left\vert  #2\right. \right]}}
\newcommand\cExpect[3]{{\mathbb{E}_{#1}\left[ \left. #3 \;\right\vert #2 \right]}}
\newcommand\cbExpect[3]{{\mathbb{E}_{#1}\big[ \big. #3 \;\big\vert #2 \big]}}
\newcommand\cBExpect[3]{{\mathbb{E}_{#1}\Big[ \Big. #3 \;\Big\vert #2 \Big]}}
\newcommand\norm[1]{\| #1 \|}
\newcommand\power[1]{\set{0,1}^{#1}}
\newcommand\ceil[1]{\lceil{#1}\rceil}
\newcommand\round[1]{\left\lfloor{#1}\right\rfloor}
\newcommand\half{{1\over2}}
\newcommand\quarter{{1\over4}}
\newcommand\eigth{{1\over8}}
\newcommand\defeq{\stackrel{def}{=}}
\newcommand\skipi{{\vskip 10pt}}
\newcommand\sqrtn{ {\sqrt n} }
\newcommand\cube[1]{ { \set{-1,1}^{#1} } }
\newcommand\minimum[2] {  \min\left\{{{#1}, {#2}}\right\}   }
\newcommand\spn{{\sf Span}}
\newcommand\spnn[1]{{\sf Span}({#1})}
\newcommand\inner[2]{\langle{#1},{#2}\rangle}
\newcommand\eps{\varepsilon}
\newcommand\sgn[1]{sign({#1})}
\newcommand*\xor{\mathbin{\oplus}}
\newcommand*\bxor{\mathbin{\bigoplus}}
\newcommand*\Hell[2]{{H\left({#1}, \ {#2}\right)}}
\newcommand*\cHell[4]{{H\left({#1}\vert({#2}),\ {#3}\vert({#4}) \right)}}
\newcommand*\dimm{{\sf dim}}
\newcommand*\tp{t}
\newcommand*\R{\mathbb{R}}
\newcommand*\mult[1]{\vec{#1}}
\newcommand*\zin[2]{\mu_{{\sf in}({#1})}({#2})}
\newcommand*\zout[2]{\mu_{{\sf out}({#1})}({#2})}
\newcommand*\zinout[3]{\mu_{{\sf in}({#1}), {\sf out}({#2})}({#3})}
\renewcommand\geq{\geqslant}
\renewcommand\ge{\geqslant}
\renewcommand\leq{\leqslant}
\renewcommand\le{\leqslant}
\newcommand\lvlfn[2]{{#1}_{={#2}}}
\newcommand\approxlvlfn[2]{{#1}_{\approx{#2}}}
\newcommand{\cupdot}{\mathbin{\mathaccent\cdot\cup}}
\newcommand{\M}{\mathcal{M}}
% unused original malformed degree macros omitted
% unused original malformed degree macros omitted

\newcommand{\rom}[1]{\uppercase\expandafter{\romannumeral #1\relax}}


\makeatletter
\def\moverlay{\mathpalette\mov@rlay}
\def\mov@rlay#1#2{\leavevmode\vtop{
   \baselineskip\z@skip \lineskiplimit-\maxdimen
   \ialign{\hfil$\m@th#1##$\hfil\cr#2\crcr}}}
\newcommand{\charfusion}[3][\mathord]{
    #1{\ifx#1\mathop\vphantom{#2}\fi
        \mathpalette\mov@rlay{#2\cr#3}
      }
    \ifx#1\mathop\expandafter\displaylimits\fi}
\makeatother
\newcommand{\bigcupdot}{\charfusion[\mathop]{\bigcup}{\cdot}}


\newcommand{\ZZ}{\mathbb{Z}}
\newcommand{\RR}{\mathbb{R}}
\newcommand{\ga}{\gamma}
\newcommand{\la}{\lambda}
\newcommand{\al}{\alpha}
\newcommand{\si}{\sigma}
\newcommand{\De}{\Delta}
\renewcommand{\r}{\hat{r}}
\newcommand{\ve}{\varepsilon}
\newcommand{\ra}{\rightarrow}

\def\gg{\gtrsim}
\def\ll{\lesssim}
\def\ee{\asymp}
\title{Permutation-symmetric lattice tiling bodies require surface area at least n over square root log n}
\author{Mark Braverman, Dor Minzer}
\date{doi:10.4230/LIPIcs.CCC.2021.5; CC BY 4.0}
\begin{document}
\maketitle
\noindent\textit{Editorial note.} Public source contains only corresponding substantive formally published CC BY material; nonexclusive native used solely as transcription aid and not publicly referenced. Blank excluded ranges preserve original locator lines. The wrapper preserves formal global mainTheorem5/6 labels and original chapter offsets; selected source-local reference closure is complete. Missing bbm font indicator rendered as ordinary bold1, and two unused malformed indeg/outdeg macros omitted; author formulas and hypotheses otherwise unchanged. All production/validation solely onYupeng viaSSH.\par\medskip
\noindent\textit{Editorial note.} Transcription correction against officialPDFpage11: published membership j in K\_i is rendered with the in macro; nativev1 had an invalid math-mode dotless-i command. Original formal author expression retained, raw native unchanged.\par\medskip
\noindent\textit{Allowed standard prerequisites (editorial).}\begin{itemize}
\item Gaussian estimates, elementary concentration, geometry of surface area and noise sensitivity. Full original source-local geometric/concentration statements and proofs included; cited standard background not recursively reproved
\end{itemize}\medskip
\section{Introduction}

A body $D\subseteq\mathbb{R}^n$ is said to be tiling the Euclidean space $\mathbb{R}^n$, if its translations by $\mathbb{Z}^n$
cover the entire space and have disjoint interiors. The foam problem asks for the least surface area a tiling body
$D$ can have. The problem had been considered by mathematicians already in the 19th century \cite{Thomson87},
and it also appears in chemistry, physics and engineering \cite{chemistry}. More recently, the problem had received
significant attention in the theoretical computer science community due to its strong relation with the parallel repetition problem~\cite{FKO,KROW,AK}.

The simplest example for a body that tiles the Euclidean space is the solid cube, $D = [0,1]^n$, which has surface
area $2n$. At first glance, one may expect the solid cube to be the best example there is, or more modestly
that any tiling body would need to have surface area $\Omega(n)$. The main results of \cite{KROW,AK}
show that this initial intuition is completely false, and that there are far more efficient tiling bodies whose surface area is $O(\sqrt{n})$.
This is surprising, since spheres --- which are the minimizers of surface area among all bodies with a given, fixed volume
(in this case volume $1$), have $\Theta(\sqrt{n})$ surface area and seem to be very far from forming a tiling of $\mathbb{R}^n$.
As we will shortly discuss, the existence of such surprising tiling body is intimately related to the existence of another
surprising object -- namely non-trivial strategies for $2$-prover-$1$-round games, repeated in parallel. The main goal of this
paper is to understand the symmetric variant of the foam problem, which is closely related to the symmetric variant of
parallel repetition.


\setcounter{thm}{3}\setcounter{subsection}{2}\subsection{Our results}
In this paper, our main object of study mainly are tilings of $\mathbb{R}^n$ using a \emph{symmetric body}.
\begin{definition}
  A set $D\subseteq\mathbb{R}^n$ is called symmetric if
  for any $\pi\in S_n$ and $x\in\mathbb{R}^n$, it holds that
  $x\in D$ if and only if $\pi(x)\in D$.
\end{definition}

The main question we consider, is what is the least surface area a symmetric tiling body can have.
Again, one has the trivial example of the solid cube $D = [0,1]^n$, but inspired by the non-symmetric
variant of the problem, one may expect there to be better examples. We first show that while this is
possible, the savings are much milder, and can be at most a multiplicative factor of $\sqrt{\log n}$.

\begin{thm}\label{thm:lb}
	Any symmetric tiling body $D$ of volume $1$
	with piecewise smooth surface has surface area at least $\Omega\left(\frac{n}{\sqrt{\log n}}\right)$.
\end{thm}


Besides the quantitative result itself, we believe the argument used in the proof of Theorem~\ref{thm:lb}
carries with it a lot of intuition regarding the additional challenge that the symmetric variants of the foam problem
and the parallel repetition posses, and we hope that this intuition will help us to develop better understanding
of symmetric parallel repetition in general. We remark that our proof actually shows a lower bound on the ``noise sensitivity''
parameter of the body, which is known to be smaller than the surface area of the body.

We complement Theorem~\ref{thm:lb} with a randomized construction showing that $O(\sqrt{\log n})$ savings
are indeed possible.
\begin{thm}\label{thm:ub}
	There exists a symmetric tiling body $D$ of volume $1$
	with piecewise smooth surface that has surface area $O\left(\frac{n}{\sqrt{\log n}}\right)$.
\end{thm}

\setcounter{thm}{8}\section{Preliminaries}\label{sec:prelim}
\paragraph{Notations.}
We write $X\ll Y$ or $X = O(Y)$ to say that
there exists an absolute constant $C >0$ such that $X\leq C\cdot Y$,
and similarly write $X\gg Y$ or $X = \Omega(Y)$ to say that
there exists an absolute constant $c >0$ such that $X\geq c\cdot Y$.
We write $X\ee Y$ or $X = \Theta(Y)$ to say that $Y\ll X\ll Y$.


We denote random variables by boldface letters such as ${\bf x}$ and $\bm{\De}$.
We denote by $\mathcal{N}(\mu,\sigma^2)$ the distribution of a standard
Gaussian random variable with mean $\mu$ and variance $\sigma^2$, and by
$\mathcal{N}(\vec{\mu}, \Sigma)$ the distribution of a multi-dimensional
Gaussian random variable with means $\vec{\mu}$ and covariance matrix $\Sigma$.
\subsection{Needles}

\begin{definition}
  Let $\delta>0$, and let $a\in\mathbb{R}^n$. A random $\delta$-needle is a line $\ell_{a,{\bf u}} = \sett{a+t\cdot {\bf u}}{t\in[0,1]}$
  where the direction vector ${\bf u}$ is a chosen as a standard Gaussian $\mathcal{N}(0,\delta I_n)$.
\end{definition}

Given a tiling body $D$, a random $\delta$-needle from $D$ is a random $\delta$-needle $\ell_{{\bf a}, {\bf u}}$
where ${\bf a}\in D$ is chosen uniformly. Random needles are a useful tool to measure the surface area of a $D$, as shown
in the following two lemmas. First, given a tiling body $D$ and a needle $\ell_{a, u}$, we may think of the needle
as ``wrapping around'' around $D$,  i.e.\ its points are taken modulo $D$. We denote this ``wrapped around'' line by
$\tilde{\ell}_{a, u}$. We will use the following formula from \cite{Santalo}; the case $n=2$ is formula (8.10) therein,
and the extension to general $n$ is discussed in page 274.
\begin{lemma}\label{lem:santalo}
  There is a constant $C_n = \Theta(1)$, such that the following holds.
  Let $S$ be a piecewise smooth surface in a tiling body $D$ of volume $1$, and let $\delta>0$.
  Then
  \[
   \Expect{{\bf a}\in D, {\bf u}\sim\mathcal{N}(0,\delta I_n)}{\card{\tilde{\ell}_{{\bf a}, {\bf u}}\cap S}} = C_n \cdot\sqrt{\delta} \cdot {\sf area}(S).
  \]
\end{lemma}


\begin{lemma}\label{lem:area_bd}
  Let $D$ be a tiling body of volume $1$, and let $\delta>0$. Then
  \[
    \Prob{{\bf a}\in D, {\bf u}\sim\mathcal{N}(0,\delta I_n)}{\ell_{{\bf a},{\bf u}}\cap \partial D\neq \emptyset} \leq \Theta(\sqrt{\delta}) {\sf area}(\partial D).
  \]
\end{lemma}
\begin{proof}
  Set $S = \partial D$, and note that whenever $\ell_{a,\delta u}\cap \partial D\neq \emptyset$, we have that $\card{\tilde{\ell}_{a,\delta u}\cap S}\geq 1$.
  Hence by the previous lemma we get that
  \[
  \Prob{{\bf a}\in D, {\bf u}\sim\mathcal{N}(0,\delta I_n)}{\ell_{{\bf a},{\bf u}}\cap \partial D\neq \emptyset}
  \leq
  \Expect{{\bf a}\in D, {\bf u}\sim\mathcal{N}(0,\delta I_n)}{\card{\tilde{\ell}_{{\bf a},{\bf u}}\cap \partial D}}
  \leq  \Theta(\sqrt{\delta}) \cdot {\sf area}(\partial D).
  \qedhere
  \]
\end{proof}
We will use the above lemma to prove lower bounds on the surface area of a tiling body, by finding $\delta$ such that
the probability on the left hand side of Lemma~\ref{lem:area_bd} is at least $\Omega(1)$; this would imply that
${\sf area}(\partial D)\geq \Omega(1/\sqrt{\delta})$.

\subsection{Basic useful properties of tiling bodies}
\begin{lemma}\label{lem:diff}
  Let $D\subseteq \mathbb{R}^n$ be a symmetric body, such that for all $z\in\mathbb{Z}^n\setminus\set{0}$
  we have $D\cap (D+z)  = \emptyset$, and let $x\in D$. Then for every $1\leq i,j\leq n$,
  if $x_i - x_j \in \mathbb{Z}$, then $x_i = x_j$.
\end{lemma}
\begin{proof}
  Assume towards contradiction $x_i - x_j$ is a non-zero integer $k$, and let $S_{i,j}\in S_n$ be the permutation that
  maps $i$ to $j$, $j$ to $i$ and has any $r\neq i,j$ as a fixed point. Since $D$ is symmetric, we have
  that $S_{i,j}(x)\in D$. Also, we have
  \[
    x - S_{i,j}(x) = (x_i - x_j) (e_i - e_j) = k(e_i-e_j),
  \]
  where $e_i$ is the $i$th element in the standard basis. In other words, we get that $x = S_{i,j}(x) + z$ for non-zero $z\in \mathbb{Z}^n$,
  and therefore $x\in D + z$. This contradict the fact that $D$ and $D+z$ are disjoint.
\end{proof}

\begin{lemma}\label{lem:uniform}
  Let $D$ be a volume $1$ tiling body, and choose $a=(a_1,\ldots,a_n)\in D$ uniformly
  at random. Then the random variable $(a_1\hspace{-1ex}\pmod{1},\ldots,a_n\hspace{-1ex}\pmod{1})$ is uniform over $[0,1)^n$.
\end{lemma}
\begin{proof}
  Sample ${\bf x}\in[0,1)^n$, and take ${\bf a} = {\bf x} \pmod{D}$. Note that the distribution of
  ${\bf a}$ is uniform over $D$. Indeed, for that we note that the map $x\to x\pmod{D}$ is bijection from
  $[0,1)^n$ to $D$:
  otherwise, there were $x\neq x'$ in $[0,1)^n$ that are equal mod $D$, and therefore differ by non-zero lattice point
  (which is clearly impossible). Now as the distribution of ${\bf a}\pmod{1}$ is just ${\bf x}$, the claim follows.
  \end{proof}


\section{The lower bound: proof of Theorem~\ref{thm:lb}}\label{sec:lb}

In this section, we prove the lower bound on the surface area of a
symmetric tiling body $D$.
Throughout, we will have two parameters: $\sigma$, which is magnitude of each coordinates in the needle we consider
(which will be of order $\frac{\sqrt{\log n}}{n}$), and an auxiliary parameter $Z$ (which will be of order $\frac{n}{\log n}$).
Let $D$ be a symmetric tiling body containing $0$.
We denote by ${\bf a}$ a random point in $D$,
and by ${\bf u}$ a Gaussian vector $\mathcal{N}(0,\sigma^2 I_n)$. We will prove that
$\Prob{{\bf a}, {\bf u}}{\ell_{{\bf a},{\bf u}}\not\subseteq D} = \Omega(1)$,
which by Lemma~\ref{lem:area_bd} implies that ${\sf area}(\partial D)\geq \Omega(1/\sigma)$. As
$\sigma = \Theta(\sqrt{\log n}/n)$, this would establish Theorem \ref{thm:lb}.

\paragraph{Notations.}
For $x,y\in\mathbb{R}$, define
\[
d(x,y):= \min_{z\in \ZZ, z\neq 0} |(x+z)-y|\in [0,1].
\]
To gain some intuition for the definition of $d(x,y)$, suppose $x$ and $y$ are two entries of a point $a\in D$.
Clearly, if $d(x,y)$ is small, then $x,y$ nearly differ by an integer $z\neq 0$, and this says that the point $a$ is
somewhat close to the boundary of $D$ (in the sense that Lemma~\ref{lem:diff} could kick in if we move along a direction
that decreases this distance).


Our argument will indeed inspect $d(a_i,a_j)$ for all distinct $i,j\in[n]$ and the way they change along a random direction.
A key measure that we will keep track of is the energy of a point $a\in D$, defined by
\begin{equation*}
\Psi(a):= \sum_{i<j} e^{-Z \cdot d(a_i,a_j)}.
\end{equation*}
We show that for ${\bf a}\in_R D$ and ${\bf u}\sim \mathcal{N}(0,\sigma^2 I_n)$, if $\ell_{{\bf a}, {\bf u}}\subseteq D$ with
probability close to $1$, then the energy of ${\bf a}$ increases along the line $\ell_{{\bf a},{\bf u}}$ with high probability, and
in particular that $\Psi({\bf a}+{\bf u}) > \Psi({\bf a})$. We then argue that with high probability, this should be the case for
the point ${\bf a}$ as well as for ${\bf a} - {\bf u}$, hence $\Psi({\bf a}+{\bf u}) > \Psi({\bf a}-{\bf u})$ with high probability.
This event however can happen with probability at most $0.5$ by symmetry, hence completing the proof.


\subsection{Analyzing the energy along a random line}
By definition of $d(x,y)$, we either have $d(x,y) = (x+z-y)$ or $d(x,y) = -(x+z-y)$ for some $z\in\mathbb{Z}\setminus\set{0}$,
and this sign determines whether $x,y$ need to move in different directions or the same direction for $d(x,y)$ to get smaller.
To capture this, we denote
\[
\ga(x,y):=\left\{ \begin{array}{ll}
+1 & \text{ if $d(x,y)=x+z-y$ for some $z\in\ZZ$, $z\neq 0$},\\
-1 & \text{ otherwise}.	
\end{array} \right.
\]
Next, we discuss the {\em energy} of a configuration, which is the key concept used in the proof.
Let $Z$ be a parameter to be chosen later (of the order $n/\log n$).
As stated earlier, our goal is to analyze the behaviour of $\Psi(a)$ along a random $\sigma^2$-needle from $a$
in direction $u$. Towards this end, note that we expect (at least if $u_i, u_j$ are small) that $d(a_i+u_i, a_j+u_j) = d(a_i,a_j) + \ga(a_i,a_j) (u_j-u_j)$,
hence expect $\Psi(a+u)$ to be close to
\begin{equation*}
\Psi(a,u):=  \sum_{i<j} e^{-Z \cdot (d(a_i,a_j)+\ga(a_i,a_j)\cdot (u_i-u_j))}.
\end{equation*}
Indeed, this is the content of the following claim.

\begin{claim}
	\label{cl:1}
	Suppose $|u_i|\leq 1/20$ for all $i$, and $a+[0,1]\cdot u \subset D$, then
	$$
	|\Psi(a+u)-\Psi(a,u)|\leq n^2\cdot e^{-Z/4}.
	$$
\end{claim}

\begin{proof}
   We consider the contribution of each pair $(i,j)$ to $\Psi(a+u)$ and $\Psi(a,u)$ separately.
   Without loss of generality we may only consider pairs $i,j$ that $\ga(a_i,a_j)=1$,
   and thus $d(a_i,a_j)= a_i - a_j + z$ for some $z\in \ZZ$, $z\neq 0$.
   Let
   \[
   d= a_i - a_j + z + (u_i-u_j) = (a_i+u_i)- (a_j+u_j)+z.
   \]
   First, we argue that $d\geq 0$. Otherwise, since $a_i-a_j + z \geq 0$ it follows by continuity that there is
   $\lambda\in[0,1]$ such that $a_i - a_j + z + \lambda(u_i-u_j) = 0$, and hence the point $a+\lambda u$ has entries
   that differ by an integer $z\neq 0$, and this contradicts Lemma \ref{lem:diff} (as $a+\lambda u \in D$).
   We now consider two cases:
   \begin{itemize}
     \item {\bf Case 1}: $d \in [0,0.5]$. In this case, we have $d(a_i+u_i, a_j+u_j) = d$, and thus the contribution of the pair $(i,j)$ to both sums is the same ($e^{-Z\cdot d}$).
     \item {\bf Case 2}: $d > 0.5$. Since $\card{u_i-u_j}\leq 0.1$, it follows that $d(a_i,a_j) = d -(u_i-u_j) > 0.4$, which implies $d(a_i+u_i,a_j+u_j) > 0.3$. Therefore, the
     contribution to $\Psi(a,u)$ from $i,j$ is at most $e^{-0.4\cdot Z}$ and to $\Psi(a+u)$ is at most $e^{-0.3 \cdot Z}$, and in particular $(i,j)$ contributes (in absolute value)
     at most $e^{-Z/4}$ to the difference between the sums.
   \end{itemize}
   Taking a sum over all pairs $(i,j)$ concludes the proof.
\end{proof}


\subsection{Analyzing the expectation and variance of $\Psi(a,{\bf u})$}
Next,
we consider $\Psi(a,{\bf u})$ as a random variable over the choice of ${\bf u}$
and compute its expectation and variance. In both computations we will use
the well-known fact that $\E[e^{-Z\cdot N(0, c^2)}]=e^{Z^2 c^2/2}$ for all $c> 0$.

\begin{claim}
\label{cl:2}
For every $a\in\mathbb{R}^n$ we have $\Expect{{\bf u}\sim \mathcal{N}(0,\sigma^2 I_n)}{\Psi(a,{\bf u})} =\Psi(a)\cdot e^{(Z \cdot \si)^2}$.
\end{claim}

\begin{proof}
    By linearity of expectation we have that
	\[
	\Expect{{\bf u}\sim \mathcal{N}(0,\sigma^2I_n)}{\Psi(a,{\bf u})}
= \sum_{i<j} e^{-Z \cdot d(a_i,a_j)}\cdot \Expect{{\bf u}\sim \mathcal{N}(0,\sigma^2I_n)}{e^{-Z\cdot \ga(a_i,a_j)\cdot ({\bf u}_i-{\bf u}_j)}}.
	\]
    Note that the above expectation does not depend on $i,j$: for every $i,j$ the distribution of
    ${\bf u}_i - {\bf u}_j$ is $N(0,\si^2)-N(0,\si^2)\sim N(0,2\si^2)$, so it is symmetric around $0$ and thus
    the sign $\ga(a_i,a_j)$ does not affect the expectation. Hence we have
    \[
    \Expect{{\bf u}\sim \mathcal{N}(0,\sigma^2I_n)}{\Psi(a,{\bf u})}   = \Psi(a) \cdot \E [e^{Z\cdot N(0,2\sigma^2)}]
	=\Psi(a)\cdot e^{Z^2\si^2}.\qedhere
    \]
\end{proof}


Next, we turn our attention into upper bounding the variance of $\Psi(a,{\bf u})$, and
for that we first define the notion of {\em good} points $a\in D$ and prove two preliminary
claims.
We say a point $a$ is {\em good} if any interval of length $(10\log n)/n$
on the torus contains at least $\log n$ and at most $100 \log n$ coordinates from $a\hspace{-1ex}\pmod{1}$.
Note by Lemma \ref{lem:uniform}, if ${\bf a}$ is chosen randomly from $D$ then ${\bf a}\pmod{1}$ is uniform
over $[0,1)^n$ and by Chernoff bound is easily shown to be good with probability $>0.999$.

We first show that good points have high energy.
\begin{claim}\label{cl:6}
There exists $c_2>0$, such that for
$Z = 0.1 \frac{\log n}{n}$, if
$a$ is good then $\Psi(a)> c_2 \log^2 n$.
\end{claim}

\begin{proof}
 Partition the torus $[0,1)$ into $m = n/(10 \log n)$ disjoint intervals of length $1/m = (10 \log n)/n$ each.
 We say that $I_i$ is {\em unanimous}, if there is $b_i\in \RR$ (called anchor) such that
 (1) $b_i\hspace{-1ex}\pmod{1}$ is the middle of $I_i$, and (2) for the majority of points $a_j\in I_i$, $|a_j-b_i|< 1/m$.

 We consider two cases:

\noindent {\bf Case 1: There is an interval $I_i$ that is not unanimous.} Note that there are at least $\log n$ coordinates $j$ of $a$ such that $a_j\in I_i$.
Let $j^{\star}$ be such coordinate, and write $a_{j^{\star}} = z_{j^{\star}} + \{a_{j^{\star}}\}$ where $z_{j^{\star}}\in\mathbb{Z}$ and
$\{a_{j^{\star}}\}$ is the fractional part of $a_{j^{\star}}$. Consider
$b = z_{j^{\star}} + m_i$ where $m_i$ is the middle of $I_i$. Then since $I_i$ is not unanimous, $b$ is not an anchor of it and so there are at least $\half \log n$
coordinates of $a$, say $(a_k)_{k\in K_{i,j^{\star}}}$ that mod $1$ are in $I_i$, and $\card{a_k - b}\geq 1/m$. Writing $a_k = z_k + \{a_k\}$, we observe that $z_k\neq z_{j^{\star}}$,
since otherwise $\card{a_k - b}= \card{\{a_k\} - m_i} \leq 1/(2m)$. Hence the difference $a_k-a_{j^{\star}}$
is $10\log n/n$ close to an integer $z_k-z_{j^{\star}}\neq 0$, and so $d(a_k,a_{j^{\star}})\leq 10\log n/n$, and the contribution of $\Psi(a)$ is
at least $e^{-1}$. Summing we get
\[
\Psi(a)\geq
\half
\sum\limits_{j^{\star}: a_{j^{\star}} \in I_i}
\sum\limits_{k\in K_{i,j^{\star}}}{e^{-Z d(a_k,a_{j^{\star}})}}
\geq
\half
\sum\limits_{j: a_j \in I_i}
e^{-1} \card{K_{i,j^{\star}}}
\geq \frac{1}{4e} \log^2 n.
\]

\noindent {\bf Case 2: All intervals are unanimous.} Let $b_i$ be an anchor of $I_i$. Note that since the fractional part of two adjacent anchors,
i.e.\ of $b_i,b_{i+1}$, are $1/m$ apart, we have that either $\card{b_i - b_{i+1}}\leq 1/m$ or $\card{b_i-b_{i+1}}\geq 1-1/m$. We claim there exists $i$ for
which the latter condition holds. To see this, assume that for all $i=1,\ldots,m-1$ we have that the first condition holds. Then we have
$b_i = z + i \frac{10\log n}{n}$ for some $z\in \mathbb{Z}$ for all $i=1,\ldots,m$, and hence $\card{b_m - b_1}\geq 1-1/m$ (and the condition holds for $i=m$).

Thus, we fix $i$ such that $\card{b_i-b_{i+1}}\geq 1-1/m$, and thus $b_i - b_{i+1} = z + \alpha$ for $z\neq 0$ and $\card{\alpha}\leq 1/m$.
Let $K_i$ be the coordinates $j$ of $a$ such that $\card{a_j - b_i}\leq 1/m$ for $j\in K_i$
and similarly define $K_{i+1}$.
We have that $a_r - a_j = z + \alpha + (a_r - b_{i+1}) + (a_j - b_i)$, hence $a_r - a_j = z + \beta$ for $\card{\beta}\leq 3/m$ for all $r\in K_{i+1}$, $j\in K_i$.
Thus $d(a_r,a_j)\leq 3/m$, and we get
\[
\Psi(a)\geq \card{K_i}\card{K_{i+1}} e^{-Z\cdot 3/m}\geq \frac{1}{4} e^{-3} \log^2 n\qedhere
\]
\end{proof}


Let $C_i = \sum_{j\neq i} e^{-Z\cdot d(a_i,a_j)}$
be the contribution of $a_i$ to $\Psi(a)$. Note that $\Psi(a)= \frac{1}{2}\sum_i C_i$.

\begin{claim}\label{cl:7}
There exists $c_3>0$, such that if $a$ is good, then for all $i$ we have
$C_i < c_3 \Psi(a)/\log n$.
\end{claim}
\begin{proof}
	Note that $d(a_i,a_j)\ge | \{a_i\} - \{a_j\}|$. Since any interval of length $10\log n/n$ on the torus contains at most $100\log n$ points of $a$,
    we have that the number of $j$'s such that $|\{a_i\} - \{a_j\}|$ is between $10 \log n/n \cdot k$ and $10 \log n/n \cdot (k+1)$ is at most $200\log n$
    (for all $k$). Therefore,
	\[
	C_i < 200 \log n\cdot \sum_{k=0}^{\infty} e^{-Z \cdot k \cdot (10 \log n )/n} = 200 \log n \cdot \sum_{k=0}^{\infty}e^{-k}
    \leq 400 \log n.
	\]
Using Claim~\ref{cl:6}, we may bound $\log n\leq \frac{1}{c_2}\frac{\Psi(a)}{\log n}$, finishing the proof.
\end{proof}


We are now ready to bound the variance of $\Psi(a,{\bf u} )$.
\begin{claim}\label{cl:3}
There exists $c_1>0$ such that the following holds.
Let $Z= n/10\log n$, let $a\in \mathbb{R}^n$ be good and let ${\bf u} \sim \mathcal{N}(0,\sigma^2 I_n)$.
Then
\[
	{\sf var}_{\bf u} [\Psi(a,{\bf u} )] \leq \frac{c_1}{\log n} \cdot( e^{4(Z \cdot \si)^2}
	-e^{2(Z \cdot \si)^2})\cdot \Psi(a)^2.
\]
\end{claim}

\begin{proof}
Using Claim \ref{cl:2} to compute the expectation of $\Psi(a,{\bf u})$, we have by definition that
\begin{align*}
&{\sf var}_{\bf u}  ( \Psi(a,{\bf u} ))
= \E_{\bf u}  \left[\left(\sum_{i<j} e^{-Z\cdot d(a_i,a_j)}\cdot (e^{Z\cdot\gamma(a_i,a_j)\cdot({\bf u} _i-{\bf u} _j)}-e^{(Z\cdot \si)^2}) \right)^2\right]\\
& =\sum_{i<j} e^{-2Z\cdot d(a_i,a_j)} \cdot\Expect{{\bf u}}{\left(e^{Z\cdot\gamma(a_i,a_j)\cdot({\bf u}_i-{\bf u}_j)}-e^{(Z\cdot \si)^2}\right)^2}\\
& + \sum_{\substack{(i,j,k) \\ \text{distinct}}}e^{-Z\cdot (d(a_i,a_j)+d(a_i,a_k))}\cdot
\Expect{{\bf u}}{(e^{Z\cdot\gamma(a_i,a_j)\cdot({\bf u}_i-{\bf u}_j)}-e^{(Z\cdot \si)^2})(e^{Z\cdot\gamma(a_i,a_k)\cdot({\bf u}_i-{\bf u}_k)}-e^{(Z\cdot \si)^2})}.
\end{align*}
Here, we used that fact that if $i,j,k,r$ are all distinct then $e^{Z\cdot\gamma(a_i,a_j)\cdot({\bf u}_i-{\bf u}_j)}$, $e^{Z\cdot\gamma(a_k,a_r)\cdot({\bf u}_k-{\bf u}_r)}$
are independent with expectation $e^{(Z\cdot \sigma)^2}$, hence the contribution of these terms is $0$. Computing, we see that
\[
\Expect{{\bf u}}{\left(e^{Z\cdot\gamma(a_i,a_j)\cdot({\bf u}_i-{\bf u}_j)}-e^{(Z\cdot \si)^2}\right)^2}
=\E\left[e^{Z\cdot N(0,8\sigma^2)}\right]-e^{2(Z\cdot \si)^2}
=e^{4(Z\cdot \si)^2}-e^{2(Z\cdot \si)^2},
\]
and
\begin{align*}
&\Expect{{\bf u}}{(e^{Z\cdot\gamma(a_i,a_j)\cdot({\bf u}_i-{\bf u}_j)}-e^{(Z\cdot \si)^2})
 (e^{Z\cdot\gamma(a_i,a_k)\cdot({\bf u}_i-{\bf u}_k)}-e^{(Z\cdot \si)^2})}\\
 &=
 \E\left[e^{(\gamma(a_i,a_j) + \gamma(a_i,a_k))Z\cdot N(0,\sigma^2)}\right]
 \E\left[e^{Z\cdot N(0,2\sigma^2)}\right]
 -e^{2(Z\cdot \si)^2}\\
 &\leq
 \E\left[e^{2Z\cdot N(0,\sigma^2)}\right]
 \E\left[e^{Z\cdot N(0,2\sigma^2)}\right]
 -e^{2(Z\cdot \si)^2}\\
 &= e^{3(Z\cdot \sigma)^2}-e^{2(Z\cdot \si)^2}.
\end{align*}
Thus, we get that
\begin{align*}
  {\sf var}_{\bf u} ( \Psi(a,{\bf u}))
  &\leq
 \sum_{i<j} e^{-2Z\cdot d(a_i,a_j)} (e^{4 (Z\cdot \si)^2}-e^{2 (Z\cdot \si)^2}) +
 \sum_{\text{$(i,j,k)$ distinct}} e^{-Z\cdot (d(a_i,a_j)+d(a_i,a_k))}(e^{3(Z\cdot \sigma)^2}-e^{2(Z\cdot \si)^2}) \\
 & \leq
 (e^{4 (Z\cdot \si)^2}-e^{2 (Z\cdot \si)^2})
 \sum_{i}
 \left(\sum_{j\neq i} e^{-2Z\cdot d(a_i,a_j)} + \sum_{j,k\neq i}
e^{-Z\cdot (d(a_i,a_j)+d(a_i,a_k))}\right)\\
&=(e^{4 (Z\cdot \si)^2}-e^{2 (Z\cdot \si)^2})\sum\limits_{i}\left(\sum_{j\neq i} e^{-2Z\cdot d(a_i,a_j)}\right)^2\\
&=  (e^{4 (Z\cdot \si)^2}-e^{2 (Z\cdot \si)^2}) \cdot \sum_i C_i^2.
\end{align*}
Therefore, using Claim~\ref{cl:7} we conclude that
\begin{multline*}
{\sf var}_{\bf u} ( \Psi(a,{\bf u}))
\leq
( e^{4 (Z\cdot \si)^2}-e^{2 (Z\cdot \si)^2}) \frac{c_3\Psi(a)}{\log n}\cdot \sum_i C_i =
\frac{2 c_3}{\log n}\cdot
( e^{4 (Z\cdot \si)^2}-e^{2 (Z\cdot \si)^2}) \cdot \Psi(a)^2.
\end{multline*}
Setting $c_1:=2c_3$ completes the proof.
\end{proof}

Putting the last two claims together, we have:

\begin{claim}
	\label{cl:4}
Let $\sigma = 10^4\sqrt{c_1} \frac{\sqrt{\log n}}{n}$ and let $a\in\mathbb{R}^n$ be good.
Then
\[
\Pr_{{\bf u}}[\Psi(a,{\bf u})>\Psi(a)+\frac{(Z\si)^4}{2}\Psi(a)]\geq 0.96.
\]
\end{claim}
\begin{proof}
	We upper bound the probability of the complement event. Using Claim \ref{cl:2} (and $e^{t}\geq 1+t+t^2/2$ for $t\geq 0$), we get
	\[
	\E_{\bf u}[\Psi(a,{\bf u})]\geq \Psi(a)\cdot \left(1+(Z \si)^2+\frac{(Z\si)^4}{2}\right).
	\]
	Hence
	\[
	\Pr_{{\bf u}}\left[\Psi(a,{\bf u})\leq \Psi(a)+\frac{(Z\si)^4}{2}\Psi(a)\right] \leq \Pr_{{\bf u}}\left[\card{\Psi(a,{\bf u})- \E_{{\bf u'}}[\Psi(a,{\bf u'})]}\geq
\Psi(a)\cdot  (Z \si)^2\right].
	\]

    We want to upper bound the probability of the last event using Chebyshev's inequality.
    Since $a$ is good, the conclusion of Claim~\ref{cl:3} holds.
    Since $Z \si =o(1)$, for large enough $n$ we get
	\[
	{\sf var}_{\bf u} [\Psi(a,{\bf u})]
    \leq
    \frac{c_1}{\log n} \cdot( e^{4(Z \cdot \si)^2}
	-e^{2(Z \cdot \si)^2})\cdot \Psi(a)^2
    \leq \frac{c_1}{\log n} \cdot \Psi(a)^2\cdot 8(Z \si)^2.
	\]
    Therefore, applying Chebyshev's inequality we see the probability in question is at most
    \[
	\frac{{\sf var}_{\bf u} [\Psi(a,{\bf u})]}{\Psi(a)^2 \cdot  (Z \si)^4} \leq
	\frac{c_1 \cdot \Psi(a)^2 \cdot 8 (Z \si)^2}{(\log n)\cdot \Psi(a)^2 \cdot  (Z \si)^4 }
    = \frac{8c_1}{(\log n)\cdot (Z\sigma)^2 } = \frac{4 c_1}{10^2 c_1} = 0.04.\qedhere
    \]
\end{proof}


\subsection{Finishing the argument}
For each $u$, denote $\eps_u = \Prob{{\bf a}\in D}{\ell_{{\bf a},u}\not\subseteq D}$,
and denote $\eps = \Expect{{\bf u}\sim\mathcal{N}(0,\sigma^2 I_n)}{\eps_{\bf u}} = \Prob{{\bf a}, {\bf u}}{\ell_{{\bf a}, {\bf u}}\not\subseteq D}$.
\begin{claim}\label{cl:tv5}
For each $u$, $\mathcal{D}_{TV}[{\bf a}; {\bf a}-u]\leq \eps_u + \eps_{-u}$.
\end{claim}
\begin{proof}
  Let $K$ be a Borel set.
  Note that it is enough to show that
  (1) if $K\subseteq D$ then $0\leq \Prob{{\bf a}\in D}{{\bf a}\in K} - \Prob{{\bf a}\in D}{{\bf a}-u\in K}\leq \eps_u$,
  and (2) if $K\subseteq\bar{D}$, then $-\eps_{-u}\leq \Prob{{\bf a}\in D}{{\bf a}\in K} - \Prob{{\bf a}\in D}{{\bf a}-u\in K}\leq 0$.
  Indeed, given both (1) and (2), the triangle inequality implies for any Borel set $K\subseteq\mathbb{R}^n$,
  \begin{align*}
  &\card{\Prob{{\bf a}\in D}{{\bf a}\in K} - \Prob{{\bf a}\in D}{{\bf a}-u\in K}}\\
  &\leq
  \card{
  \Prob{{\bf a}\in D}{{\bf a}\in K\cap D}-\Prob{{\bf a}\in D}{{\bf a}-u\in K\cap D}
  +\Prob{{\bf a}\in D}{{\bf a}\in K\setminus D}-\Prob{{\bf a}\in D}{{\bf a}-u\in K\setminus D}}
  \leq \eps_u+\eps_{-u}.
  \end{align*}

  To prove (1), note that
  $\Prob{{\bf a}\in D}{{\bf a}\in K}=\mu(K)$
  and
  \[
  \Prob{{\bf a}\in D}{{\bf a}-u\in K} = \Prob{{\bf a}\in D}{{\bf a}\in K+u} = \mu((K+u)\cap D).
  \]
  This is at most $\mu(K+u) = \mu(K)$ (hence the expression in (1) is non-negative)
  and at least $\geq \mu(K+u) - \mu((K+u)\setminus D)= \mu(K) - \mu(K\setminus (D-u))$.
  Therefore
  \[
  0\leq \Prob{{\bf a}\in D}{{\bf a}\in K} - \Prob{{\bf a}\in D}{{\bf a}-u\in K}
  \leq \mu(K\setminus (D-u))
  \leq \mu(D\setminus (D-u))
  = \Prob{{\bf a}\in D}{{\bf a}+u\not\in D}
  \leq \eps_u.
  \]

  To prove (2), note that $\Prob{{\bf a}\in D}{{\bf a}\in K} = 0$ (hence the expression in (2) is non-positive)
  and
  \[
  \Prob{{\bf a}\in D}{{\bf a}-u\in K} \leq \Prob{{\bf a}\in D}{{\bf a}-u\not\in D} \leq \eps_{-u}.\qedhere
  \]
\end{proof}

\begin{claim}
	\label{cl:5}
	$\eps\geq 0.1$.
\end{claim}
\begin{proof}
	Let $E_1$ be the event that ${\bf a} + {\bf u}[0,1]\subseteq D$, let
    $E_2$ be the event that $\Psi({\bf a})\leq 1$, let $E_3$ be the event
	that $|{\bf u}_i|>1/20$ for some $i$ and let $E_4$ be the event that
	$\Psi({\bf a},{\bf u})>\Psi({\bf a})+\frac{(Z\si)^4}{2}\Psi({\bf a})$. Finally, let $E_5$ be the event that
	$\Psi({\bf a}+{\bf u})>\Psi({\bf a})$ and denote $E({\bf a},{\bf u}) = E_1\cap (\neg E_2)\cap (\neg E_3)\cap E_4$.
    Note that if the event $E$ holds for $a,u$, then $E_5$ also holds, since by Claim~\ref{cl:1}:
	\[
	\Psi(a+u)\geq \Psi(a,u)-n^2\cdot e^{-Z/4}> \Psi(a)+\frac{(Z\si)^4}{2}\Psi(a)-n^2\cdot e^{-Z/4}\geq \Psi(a).
	\]



    By Claim \ref{cl:6} the probability of $E_2$ is at most the probability ${\bf a}$ is bad, hence it is at most $0.005$.
    By definition, the probability of $E_1$ is $1-\eps$.
    By the union bound and Chernoff inequality, the probability of $E_3$ is $o(1)$.
	Thus, by Claim~\ref{cl:4} we have
	\begin{equation}
	\Pr_{{\bf a},{\bf u}}[E({\bf a},{\bf u})]\geq 0.96 -\eps- 0.005 - o(1)\geq 0.95-\eps.
	\end{equation}
    Fix $u$. Using Claim~\ref{cl:tv5} we get that
    \[
    \Prob{{\bf a}}{E({\bf a}-u,u)} \geq
    \Prob{{\bf a}}{E({\bf a},u)} - \mathcal{D}_{TV}[{\bf a}; {\bf a}-u]
    \geq \Prob{{\bf a}}{E({\bf a},u)} - \eps_u-\eps_{-u}.
    \]
    By the union bound, we now conclude that
    \[
    \Prob{{\bf a}}{E({\bf a}-u,u)\cap E({\bf a},u)}\geq
    1 -  \Prob{{\bf a}}{\overline{E({\bf a}-u,u)}}
    -  \Prob{{\bf a}}{\overline{E({\bf a},u)}}
    \geq 2 \Prob{{\bf a}}{E({\bf a},u)} -1-\eps_u-\eps_{-u}.
    \]
    Taking expectation over ${\bf u}$, we get that
    \[
      \Prob{{\bf a},{\bf u}}{E({\bf a}-{\bf u},{\bf u})\cap E({\bf a},{\bf u})}\geq
     2\Prob{{\bf a},{\bf u}}{E({\bf a},{\bf u})} -1-2\Expect{{\bf u}}{\eps_{\bf u}}
     \geq 0.9 - 4\eps.
    \]
    Next, when both $E(a-u,u)$ and $E(a,u)$ hold, we have by the previous observation that $E_5$ holds for both pairs $(a-u,u)$ and $(a,u)$,
    and so $\Psi(a+u) > \Psi (a) = \Psi((a-u) + u) > \Psi(a-u)$. Thus, we get that
    $\Prob{{\bf a},{\bf u}}{\Psi({\bf a}+{\bf u}) > \Psi({\bf a}-{\bf u})}\geq 0.9 - 4\eps$. On the other hand, the probability on the left hand side is at most $0.5$;
    this follows as
    $\Prob{{\bf a},{\bf u}}{\Psi({\bf a}+{\bf u}) > \Psi({\bf a}-{\bf u})}
    = \Prob{{\bf a},{\bf u}}{\Psi({\bf a}-{\bf u}) > \Psi({\bf a}+{\bf u})}$ (since the distributions of ${\bf u}$ and $-{\bf u}$ are identical) and their sum is at most $1$.
    Combining the two inequalities we get that $\eps \geq 0.1$.
\end{proof}

\begin{thebibliography}{10}

\bibitem{perfect}
Perfect sets are uncountable.
\newblock https://mathcs.org/analysis/reals/topo/proofs/pfctuncb.html.
\newblock [Online; accessed 15-April-2020].

\bibitem{AK}
Noga Alon and Bo'az Klartag.
\newblock Economical toric spines via cheeger's inequality.
\newblock {\em Journal of Topology and Analysis}, 1(02):101--111, 2009.

\bibitem{AKKST}
Sanjeev Arora, Subhash Khot, Alexandra Kolla, David Steurer, Madhur Tulsiani,
  and Nisheeth~K. Vishnoi.
\newblock Unique games on expanding constraint graphs are easy: extended
  abstract.
\newblock In {\em Proceedings of the 40th Annual {ACM} Symposium on Theory of
  Computing, Victoria, British Columbia, Canada, May 17-20, 2008}, pages
  21--28, 2008.

\bibitem{ALMSS}
Sanjeev Arora, Carsten Lund, Rajeev Motwani, Madhu Sudan, and Mario Szegedy.
\newblock Proof verification and the hardness of approximation problems.
\newblock {\em J. {ACM}}, 45(3):501--555, 1998.

\bibitem{AS}
Sanjeev Arora and Shmuel Safra.
\newblock Probabilistic checking of proofs: {A} new characterization of {NP}.
\newblock {\em J. {ACM}}, 45(1):70--122, 1998.

\bibitem{BHHRRS}
Boaz Barak, Moritz Hardt, Ishay Haviv, Anup Rao, Oded Regev, and David Steurer.
\newblock Rounding parallel repetitions of unique games.
\newblock In {\em 49th Annual {IEEE} Symposium on Foundations of Computer
  Science, {FOCS} 2008, October 25-28, 2008, Philadelphia, PA, {USA}}, pages
  374--383, 2008.

\bibitem{BKS}
Boaz Barak, Pravesh~K. Kothari, and David Steurer.
\newblock Small-set expansion in shortcode graph and the 2-to-2 conjecture.
\newblock In {\em 10th Innovations in Theoretical Computer Science Conference,
  {ITCS} 2019, January 10-12, 2019, San Diego, California, {USA}}, pages
  9:1--9:12, 2019.

\bibitem{BRRRS}
Boaz Barak, Anup Rao, Ran Raz, Ricky Rosen, and Ronen Shaltiel.
\newblock Strong parallel repetition theorem for free projection games.
\newblock In {\em Approximation, Randomization, and Combinatorial Optimization.
  Algorithms and Techniques, 12th International Workshop, {APPROX} 2009, and
  13th International Workshop, {RANDOM} 2009, Berkeley, CA, USA, August 21-23,
  2009. Proceedings}, pages 352--365, 2009.

\bibitem{BSVV}
Amey Bhangale, Ramprasad Saptharishi, Girish Varma, and Rakesh Venkat.
\newblock On fortification of projection games.
\newblock In {\em Approximation, Randomization, and Combinatorial Optimization.
  Algorithms and Techniques, {APPROX/RANDOM} 2015, August 24-26, 2015,
  Princeton, NJ, {USA}}, pages 497--511, 2015.

\bibitem{BG}
Mark Braverman and Ankit Garg.
\newblock Small value parallel repetition for general games.
\newblock In {\em Proceedings of the Forty-Seventh Annual {ACM} on Symposium on
  Theory of Computing, {STOC} 2015, Portland, OR, USA, June 14-17, 2015}, pages
  335--340, 2015.

\bibitem{DKKMS2}
Irit Dinur, Subhash Khot, Guy Kindler, Dor Minzer, and Muli Safra.
\newblock On non-optimally expanding sets in grassmann graphs.
\newblock In {\em Proceedings of the 50th Annual {ACM} {SIGACT} Symposium on
  Theory of Computing, {STOC} 2018, Los Angeles, CA, USA, June 25-29, 2018},
  pages 940--951, 2018.

\bibitem{DKKMS1}
Irit Dinur, Subhash Khot, Guy Kindler, Dor Minzer, and Muli Safra.
\newblock Towards a proof of the 2-to-1 games conjecture?
\newblock In {\em Proceedings of the 50th Annual {ACM} {SIGACT} Symposium on
  Theory of Computing, {STOC} 2018, Los Angeles, CA, USA, June 25-29, 2018},
  pages 376--389, 2018.

\bibitem{DS}
Irit Dinur and David Steurer.
\newblock Analytical approach to parallel repetition.
\newblock In {\em Symposium on Theory of Computing, {STOC} 2014, New York, NY,
  USA, May 31 - June 03, 2014}, pages 624--633, 2014.

\bibitem{FGLSS}
Uriel Feige, Shafi Goldwasser, Laszlo Lov\'{a}sz, Shmuel Safra, and Mario
  Szegedy.
\newblock Interactive proofs and the hardness of approximating cliques.
\newblock {\em J. ACM}, 43(2):268--292, March 1996.

\bibitem{FKO}
Uriel Feige, Guy Kindler, and Ryan O'Donnell.
\newblock Understanding parallel repetition requires understanding foams.
\newblock In {\em Twenty-Second Annual IEEE Conference on Computational
  Complexity (CCC'07)}, pages 179--192. IEEE, 2007.

\bibitem{FRS}
Lance Fortnow, John Rompel, and Michael Sipser.
\newblock On the power of multi-prover interactive protocols.
\newblock {\em Theor. Comput. Sci.}, 134(2):545--557, 1994.

\bibitem{GW}
Michel~X. Goemans and David~P. Williamson.
\newblock Improved approximation algorithms for maximum cut and satisfiability
  problems using semidefinite programming.
\newblock {\em J. {ACM}}, 42(6):1115--1145, 1995.

\bibitem{Holenstein}
Thomas Holenstein.
\newblock Parallel repetition: Simplification and the no-signaling case.
\newblock {\em Theory of Computing}, 5(1):141--172, 2009.

\bibitem{Khot}
Subhash Khot.
\newblock On the power of unique 2-prover 1-round games.
\newblock In {\em Proceedings of the 17th Annual {IEEE} Conference on
  Computational Complexity, Montr{\'{e}}al, Qu{\'{e}}bec, Canada, May 21-24,
  2002}, page~25, 2002.

\bibitem{KhotICM-survey}
Subhash Khot.
\newblock Inapproximability of {NP}-complete problems, discrete fourier
  analysis, and geometry.
\newblock In {\em Proceedings of the International Congress of Mathematicians
  2010}, pages 2676--2697, 2010.

\bibitem{KKMO}
Subhash Khot, Guy Kindler, Elchanan Mossel, and Ryan O'Donnell.
\newblock Optimal inapproximability results for {MAX-CUT} and other 2-variable
  csps?
\newblock {\em {SIAM} J. Comput.}, 37(1):319--357, 2007.

\bibitem{KMS1}
Subhash Khot, Dor Minzer, and Muli Safra.
\newblock On independent sets, 2-to-2 games, and {G}rassmann graphs.
\newblock In {\em Proceedings of the 49th Annual {ACM} {SIGACT} Symposium on
  Theory of Computing, {STOC} 2017, Montreal, QC, Canada, June 19-23, 2017},
  pages 576--589, 2017.

\bibitem{KMS2}
Subhash Khot, Dor Minzer, and Muli Safra.
\newblock Pseudorandom sets in grassmann graph have near-perfect expansion.
\newblock In {\em 59th {IEEE} Annual Symposium on Foundations of Computer
  Science, {FOCS} 2018, Paris, France, October 7-9, 2018}, pages 592--601,
  2018.

\bibitem{KROW}
Guy Kindler, Anup Rao, Ryan O'Donnell, and Avi Wigderson.
\newblock Spherical cubes: optimal foams from computational hardness
  amplification.
\newblock {\em Commun. {ACM}}, 55(10):90--97, 2012.

\bibitem{Moshkovitz}
Dana Moshkovitz.
\newblock Parallel repetition from fortification.
\newblock In {\em 55th {IEEE} Annual Symposium on Foundations of Computer
  Science, {FOCS} 2014, Philadelphia, PA, USA, October 18-21, 2014}, pages
  414--423, 2014.

\bibitem{Rao}
Anup Rao.
\newblock Parallel repetition in projection games and a concentration bound.
\newblock {\em SIAM Journal on Computing}, 40(6):1871--1891, 2011.

\bibitem{Razrep}
Ran Raz.
\newblock A parallel repetition theorem.
\newblock {\em {SIAM} J. Comput.}, 27(3):763--803, 1998.

\bibitem{Raz}
Ran Raz.
\newblock A counterexample to strong parallel repetition.
\newblock {\em SIAM Journal on Computing}, 40(3):771--777, 2011.

\bibitem{ros2001isoperimetric}
Antonio Ros.
\newblock The isoperimetric problem.
\newblock {\em Global Theory of Minimal Surfaces. Clay Math. Proc}, vol.
  2:175--209, 2005.

\bibitem{chemistry}
Jean-Fran{\c{c}}ois Sadoc and Nicolas Rivier.
\newblock {\em Foams and emulsions}, volume 354.
\newblock Springer Science \& Business Media, 2013.

\bibitem{SS}
Shmuel Safra and Oded Schwartz.
\newblock On parallel-repetition, {U}nique-{G}ames and {M}ax-{C}ut.
\newblock 2007.

\bibitem{Santalo}
Luis Antonio~Santal{\'o} Sors and Luis~A Santal{\'o}.
\newblock {\em Integral geometry and geometric probability}.
\newblock Cambridge university press, 2004.

\bibitem{Thomson87}
William Thomson.
\newblock On the division of space with minimum partitional area.
\newblock {\em Acta Math.}, 11:121--134, 1887.

\bibitem{TrevisanUGCSurvey}
Luca Trevisan.
\newblock On {K}hot's unique games conjecture.
\newblock {\em Bull. Amer. Math. Soc. (N.S.)}, 49(1):91--111, 2012.

\bibitem{tulsiani2014optimal}
Madhur Tulsiani, John Wright, and Yuan Zhou.
\newblock Optimal strong parallel repetition for projection games on low
  threshold rank graphs.
\newblock In {\em International Colloquium on Automata, Languages, and
  Programming}, pages 1003--1014. Springer, 2014.

\end{thebibliography}
\end{document}
